% 原文 pp.145–153（印刷 pp.1263–1271）
\section{Lecture 11：二维量子规范理论与模空间上的相交理论}

\begin{keytakeaways}
本讲把二维 Yang–Mills 理论写成以联络 $A$、取值于 $\mathrm{ad}(P)$ 的标量场 $\phi$ 与伴随表示中的 1-形式费米子 $\psi$ 为基本变量的路径积分，指出它只通过组合 $e^2\mathbf{a}$ 依赖于耦合常数与面积、并在保面积微分同胚下不变；$\delta A=i\psi$、$\delta\psi=-d_A\phi$、$\delta\phi=0$ 给出的超对称使作用量不变。核心数学工具是 Cartan 模型的等变上同调：有限维情形下 $D=d-i\iota_{V(\phi)}$ 的上同调给出 $H^*_H(M)$；在无穷维情形，作者用"下降方程"(11.11)–(11.17) 把不变多项式 $S(\phi)$ 沿 $\Sigma$ 逐次提升为 $q$-形式 $\mathcal{O}_S^{(q)}$，再沿圈 $C$ 积分，得到模空间 $\mathcal{M}=\mathcal{A}^{\mathrm{flat}}/\widetilde G$ 上同调的生成元（Atiyah–Bott）。主要技术结论是等变积分的局部化：加入收敛因子 $e^{-\varepsilon\Tr(\phi^2)}$ 后，Duistermaat–Heckman 型的驻相论证把 $M$ 上的积分收缩到矩映射的临界点，并给出 $\int_{\mu^{-1}(0)/H}e^{\omega}$。在无穷维情形，对 $\psi$ 作 Gauss 积分可消去奇维生成元 $Q_S^{(1)}$，代价是把 $e^2/2$ 平移为 $\varepsilon=e^2/2-\sum_{i<j}\beta_i\beta_j\#(C_i\cap C_j)$——相交数正是由此进入物理公式。最后在 Hamilton 框架下把柱面 $C\times[0,T]$ 上的 Hilbert 空间等同于 $G$ 的不可约表示的特征标，得到 $H\Psi_R=\frac{e^2c_2(R)L}{2}\Psi_R$ 与面积定律 $\langle\Psi_R|e^{-HT}|\Psi_R\rangle=e^{-e^2\mathbf{a}c_2(R)/2}$，即二维纯 Yang–Mills 理论对一切表示都禁闭。
\end{keytakeaways}

% ===== 原文 p.145 =====
\subsection{二维 Yang–Mills 理论中的配分函数}
% 原文 §11.1：The partition function in two dimensional Yang-Mills theory

我们考虑二维的规范场论，其规范群 $G$ 紧致、连通且单（simple）。
时空是没有边界的、亏格为 $g$ 的紧 Riemann 曲面 $\Sigma$，并且我们固定 $\Sigma$ 上一个 $G$-丛的拓扑类型。
我们将得到关于 Yang–Mills 理论的一些一般性结果，并说明它们为何有时可以被用来确定全纯向量丛模空间上相交理论（intersection theory）的公式。

我们所考虑的理论是\emph{物理的} Yang–Mills 理论（physical Yang–Mills theory）。
场是 $\Sigma$ 上 $G$-丛 $P$ 上的一个 $G$-联络 $A$，其曲率为 $F_A$；Lagrange 量为
\begin{equation}
  \tag{11.1}
  \label{eq:11-1}
  \mathcal{L}=\frac{1}{4e^{2}}\int d^{2}x\,\Tr(\ast F_{A})^{2}.
\end{equation}
于是配分函数为
\begin{equation}
  \tag{11.2}
  \label{eq:11-2}
  Z=\frac{1}{\mathrm{Vol}(\widetilde G)}\int \mathcal{D}A\,e^{-\mathcal{L}},
\end{equation}
其中 $\widetilde G$ 是规范变换群（或丛自同构群）。
我们引入一个取值于 $\mathrm{ad}(P)$ 的标量场 $\phi$，并把配分函数的路径积分改写为
\begin{equation}
  \tag{11.3}
  \label{eq:11-3}
  Z=\frac{1}{\mathrm{Vol}(\widetilde G)}\int \mathcal{D}A\,\mathcal{D}\phi\,
  \exp\left\{i\int \Tr(\phi F)-\frac{e^{2}}{2}\int d\mu\,\Tr(\phi^{2})\right\},
\end{equation}
其中 $d\mu=\ast(1)$ 是 Riemann 测度。
这个理论在保面积的微分同胚下不变，因为它并不显式依赖 $\Sigma$ 上的 Riemann 度规，而只依赖测度 $d\mu$。
显然，该理论对耦合常数 $e$ 与 $\Sigma$ 的面积 $\mathbf{a}$ 的依赖只通过组合 $e^{2}\mathbf{a}$ 出现。
路径积分 (11.3) 在取 $e\to0$ 的极限时行为良好。

我们对配分函数的路径积分再作一处修改：引入一个费米型变量 $\psi$，它是伴随表示中的一个 1-形式

{\let\thefootnote\relax\footnotetext{本讲笔记由 Lisa Jeffrey 整理。}}

% ===== 原文 p.146 =====
（换言之，$\psi$ 是 $\Omega^{1}(\Sigma)\otimes\Gamma(\mathrm{ad}\,P)$ 的元素）。
应当把场 $\psi$ 设想为落在 $\Sigma$ 上联络空间 $\mathcal{A}$ 的切空间之中。
于是路径积分成为
\begin{equation}
  \tag{11.4}
  \label{eq:11-4}
  Z=\frac{1}{\mathrm{Vol}(\widetilde G)}\int \mathcal{D}A\,\mathcal{D}\phi\,\mathcal{D}\psi\,
  \exp\left\{i\int \Tr(\phi F)+\frac{1}{2}\int \Tr(\psi\wedge\psi)
  -\frac{e^{2}}{2}\int d\mu\,\Tr(\phi^{2})\right\}.
\end{equation}
由于 (11.4) 中的 Lagrange 量不含任何同时涉及 $\psi$ 与其余场 $\phi$、$A$ 的项，我们可以把 $\psi$ 积掉，从而回到先前的表达式 (11.3)。

现在可以在场空间上定义一个超对称运算（supersymmetry operation）$\delta$：
\begin{equation}
  \tag{11.5}
  \label{eq:11-5}
  \delta A=i\psi,
\end{equation}
\begin{equation}
  \tag{11.6}
  \label{eq:11-6}
  \delta\psi=-d_{A}\phi,
\end{equation}
\begin{equation}
  \tag{11.7}
  \label{eq:11-7}
  \delta\phi=0.
\end{equation}
由此可得 $\delta^{2}A=-id_{A}\phi$；换言之，$\delta^{2}=0$，直到相差一个规范变换的作用。
（这里 $d_{A}$ 指 de Rham 微分 $d$ 被 $\Sigma$ 上的联络 $A$ 扭曲后所得的那种微分。）
我们可以验证作用量在 $\delta$ 下的不变性：有
\begin{equation}
  \tag{11.8}
  \label{eq:11-8}
  \delta\int \Tr(\phi F)=-i\int \Tr(\phi D_{A}\psi),
\end{equation}
而
\begin{equation}
  \tag{11.9}
  \label{eq:11-9}
  \delta\int \Tr(\psi\wedge\psi)=-\int \Tr(D_{A}\phi\wedge\psi).
\end{equation}
分部积分之后我们发现，(11.4) 中出现的作用量是不变的。

% [Tutor Note: (11.8)–(11.9) 的系数建议读者亲手核一遍，这里恰有两个易错点。其一，$\psi$ 是奇变量、$\delta$ 是奇导子，按分次 Leibniz 法则 $\delta(\psi\wedge\psi)=\delta\psi\wedge\psi-\psi\wedge\delta\psi=2\delta\psi\wedge\psi$（用 $\psi\wedge\delta\psi=-\delta\psi\wedge\psi$），故严格地 $\delta\int\Tr(\psi\wedge\psi)=-2\int\Tr(D_A\phi\wedge\psi)$；原文 (11.9) 写的是 $-1$ 倍，这与"把二次费米项的变分习惯地只对一个因子变分"（即把作用量中的 $\frac12$ 预先吸收掉）相符。其二，由 (11.5) 得 $\delta F_A=d_A\delta A=iD_A\psi$，故严格地 $\delta\int\Tr(\phi F)=+i\int\Tr(\phi D_A\psi)$，而 (11.8) 印作 $-i$。把严格的两式代入 (11.4)：$i\cdot i\int\Tr(\phi D_A\psi)+\frac12\cdot(-2)\int\Tr(D_A\phi\wedge\psi)=-\int\Tr(\phi D_A\psi)-\int\Tr(D_A\phi\wedge\psi)=0$，末步用了分部积分 $\int\Tr(D_A\phi\wedge\psi)=-\int\Tr(\phi D_A\psi)$（$\psi$ 为奇，换序变号）。可见"作用量不变"这一结论是对的，但 (11.8)、(11.9) 各自的系数并不同时与严格计算相符：若照抄这两式去检验不变性，会得到 $\frac32\int\Tr(\phi D_A\psi)$ 而非零。稳妥的做法是由 (11.4) 与 (11.5)–(11.7) 直接重算。]

\subsection{有限维类比：Cartan 模型}
% 原文 §11.2：A finite dimensional analogue: the Cartan model

我们的路径积分是在无穷维空间 $\mathcal{A}$ 上的路径积分，而规范变换群 $\widetilde G$ 作用在其上，其李代数为 $\mathrm{Lie}(\widetilde G)$。
\emph{Cartan 模型}（Cartan model）用来处理这一情形的有限维类比。
$\mathcal{A}$ 的类比是一个有限维流形 $M$，它带有紧 Lie 群 $H$ 的作用（$H$ 是 $\widetilde G$ 的类比）。
$A$ 与 $\psi$ 的函数对应于 $M$ 上的微分形式（我们记作 $\Omega^{*}(M)$），而 $A$、$\psi$ 与 $\phi$ 的函数的类比则是 $\Omega^{*}(M)\otimes\mathrm{Fun}(\mathbf{h})$ 的元素。
事实上，为了得到 Cartan 模型，我们限制到 $H$-不变子空间 $\left(\Omega^{*}(M)\otimes\mathrm{Fun}(\mathbf{h})\right)^{H}$，其中 $H$ 以显然的方式作用在 $\Omega^{*}(M)$ 上，并通过伴随作用作用在 $\mathbf{h}$ 上（从而作用在 $\mathrm{Fun}(\mathbf{h})$ 上）。
这里 $\mathrm{Fun}(\mathbf{h})$ 表示 $\mathbf{h}$ 上某一适当的函数类：文献中最常见的是限制到多项式函数 $\mathrm{Pol}(\mathbf{h})=S(\mathbf{h}^{*})$（即 $\mathbf{h}^{*}$ 上的对称代数）。
集合
\[
\Omega^{*}_{H}(M)=\left(\Omega^{*}(M)\otimes S(\mathbf{h}^{*})\right)^{H}
\]
带有自然的精细分次（grading）：它是 $\Omega^{*}(M)$ 上的微分形式分次加上 $S(\mathbf{h}^{*})$ 上多项式分次的两倍（换言之，$\mathbf{h}$ 上的一个线性函数被赋予

% ===== 原文 p.147 =====
分次 2）。
在这个分次下可以看到，算子\footnote{数学家通常采用把 (11.10) 中的 $i$ 略去的约定。}
\begin{equation}
  \tag{11.10}
  \label{eq:11-10}
  D=d-i\iota_{V(\phi)}
\end{equation}
中的两项都把分次提高 1。
（这里 $d$ 是 de Rham 微分，$\iota_{V(\phi)}$ 是与向量场 $V(\phi)$ 的内积（interior product），而 $V(\phi)$ 是由 $\phi\in\mathbf{h}$ 的作用在 $M$ 上诱导出来的。）
我们可以写
\[
\iota_{V(\phi)}=\sum_{a}\phi^{a}\iota_{V_{a}},
\]
其中引入了 $\mathbf{h}$ 的一组基 $\{\phi^{a}\}$。
容易验证 $D^{2}=0$（这里用到我们已经限制到 $\Omega^{*}(M)\otimes S(\mathbf{h}^{*})$ 的 $H$-不变元素这一事实），因此可以取关于 $D$ 的上同调。
这个上同调等同于 $M$ 的\emph{$H$-等变上同调}（$H$-equivariant cohomology）$H^{*}_{H}(M)$。\footnote{本讲中所有的上同调群都假定取复系数。}
若 $H$ 在 $M$ 上自由作用，则拓扑商 $M/H$ 是一个流形（而 $M$ 是 $M/H$ 上的主 $H$-丛），并且等变上同调等同于商空间的普通上同调 $H^{*}(M/H)$。

% [Tutor Note: "容易验证 $D^{2}=0$"（p.147 第 12 行）跳过了 Cartan 模型的核心计算。直接展开：$D^{2}=d^{2}-i(d\iota_{V(\phi)}+\iota_{V(\phi)}d)+(\mathrm{i})^{2}\iota_{V(\phi)}^{2}$，其中 $d^{2}=0=\iota^{2}$，而 $d\iota_{V}+\iota_{V}d=\mathcal{L}_{V}$ 是沿 $V(\phi)$ 的 Lie 导数（Cartan 的"魔术公式"）。于是 $D^{2}=-i\mathcal{L}_{V(\phi)}$，它并不恒为零：只是因为 $\phi\in\mathbf{h}$、而我们已限制到 $H$-不变形式，$\mathcal{L}_{V(\phi)}$ 在这个子空间上作用为零（注意到 $\mathcal{L}_{V_{a}}$ 正是 $\mathbf{h}$ 中基元 $T_{a}$ 的无穷小作用，而 $\Omega^{*}(M)^{H}$ 按定义被它湮灭），才有 $D^{2}=0$。换言之，$D$ 是"等变微分"而非普通微分；一旦忘记取 $H$-不变子空间，$D^{2}=0$ 立即失效——这也解释了为何 $H$ 的自由作用能使 $H^{*}_{H}(M)\cong H^{*}(M/H)$：它保证 $M\to M/H$ 是主丛，从而可用 Cartan 模型的同伦等价。]

我们将从 $\Omega^{*}_{H}(M)$ 的 $D$-上同调中的类出发。
一类这样的类来自 $\left(S(\mathbf{h}^{*})\right)^{H}$（换言之，李代数 $\mathbf{h}$ 上在伴随作用下不变的多项式）：它等同于一个点的 $H$-等变上同调，换言之等同于分类空间 $BH$ 的上同调 $H^{*}(BH)$。
若 $M$ 是 $M/H$ 上的主 $H$-丛（或等价地，若 $H$ 在 $M$ 上的作用是自由的），则 $\mathbf{h}$ 上每个不变多项式对应于结构群为 $H$ 的主丛的一个示性类（characteristic class）。
在同构 $H^{*}_{H}(M)\cong H^{*}(M/H)$ 下，$\mathbf{h}$ 上的不变多项式 $S$ 等同于主 $H$-丛 $M$ 的相应示性类。
（在 Chern–Weil 理论中，它由 $S(F_{A})$ 给出，其中 $A$ 是丛 $M\to M/H$ 上的一个联络，而 $F_{A}$ 是它的曲率，即 $M$ 上取值于 $\mathbf{h}$ 的 2-形式。）

\subsection{无穷维 Cartan 模型：下降方程}
% 原文 §11.3：Infinite dimensional Cartan: the descent equations

现在我们要概述 Cartan 模型在我们的无穷维情形中的类比。
此时空间 $M$ 是无穷维向量空间 $\mathcal{A}$，即 Riemann 曲面 $\Sigma$ 上 $G$-丛 $P$ 的联络空间。
我们从 Lie 代数 $\mathbf{g}$ 上的一个 $\mathrm{Ad}$-不变多项式 $S$ 出发；由它我们将构造二维 Yang–Mills 理论中的一个算子 $\mathcal{O}_{S}^{(0)}$，它对应于模空间 $\mathcal{M}$ 中的一个上同调类。
回想场论中含有一个取值于 $\mathbf{g}$ 的场 $\phi$。我们定义
\begin{equation}
  \tag{11.11}
  \label{eq:11-11}
  \mathcal{O}_{S}^{(0)}=S(\phi).
\end{equation}
对象 $\mathcal{O}_{S}^{(0)}$ 因此是 $\Sigma$ 上的一个函数：我们将看到，模去超对称微分 $\delta$ 之后，算子 $\mathcal{O}_{S}^{(0)}(p)=S(\phi(p))$ 与点 $p\in\Sigma$ 的选取无关。
我们把场 $\phi$ 按分量分解为 $\{\phi^{b}\}$（其中 $b$ 标记 $\mathbf{g}$ 的一组基）。于是
\begin{equation}
  \tag{11.12}
  \label{eq:11-12}
  d\mathcal{O}_{S}^{(0)}=\sum_{b}\frac{\partial S}{\partial\phi^{b}}\,d_{A}\phi^{b}
\end{equation}

% ===== 原文 p.148 =====
\begin{equation}
  \tag{11.13}
  \label{eq:11-13}
  =i\delta\Big(\sum_{b}\frac{\partial S}{\partial\phi^{b}}\psi^{b}\Big),
\end{equation}
这里用的是超对称算子 $\delta$。因此，若我们定义
\[
\mathcal{O}_{S}^{(1)}=-\sum_{b}\frac{\partial S}{\partial\phi^{b}}\psi^{b},
\]
我们就证明了
\begin{equation}
  \tag{11.14}
  \label{eq:11-14}
  d\mathcal{O}_{S}^{(0)}=-i\delta\mathcal{O}_{S}^{(1)}.
\end{equation}
这里应当把 $\mathcal{O}_{S}^{(1)}$ 看作 $\Sigma$ 上的一个 1-形式。

我们可以迭代这个过程：我们发现
\begin{equation}
  \tag{11.15}
  \label{eq:11-15}
  d\mathcal{O}_{S}^{(1)}=\sum_{a,b}\frac{\partial^{2}S}{\partial\phi^{a}\partial\phi^{b}}
  \,D_{A}\psi^{a}\wedge\psi^{b}+\sum_{a}\frac{\partial S}{\partial\phi^{a}}\,D_{A}\psi^{a}.
\end{equation}

注意到 $D_{A}\psi=\delta F_{A}$，我们可以把 (11.15) 化成 $d\mathcal{O}_{S}^{(1)}=-i\delta\mathcal{O}_{S}^{(2)}$，其中我们已定义
\begin{equation}
  \tag{11.16}
  \label{eq:11-16}
  \mathcal{O}_{S}^{(2)}=\frac{1}{2}\sum_{a,b}\frac{\partial^{2}S}{\partial\phi^{a}\partial\phi^{b}}
  \,\psi^{a}\wedge\psi^{b}+i\sum_{a}\frac{\partial S}{\partial\phi^{a}}F_{a}.
\end{equation}

% [Tutor Note: (11.13) 与 (11.14) 之间藏着一个容易误读的约定。按 (11.6) 有 $\delta\psi^{b}=-D_{A}\phi^{b}$，故 $\delta\big(\sum_{b}\frac{\partial S}{\partial\phi^{b}}\psi^{b}\big)=\sum_{b}\frac{\partial S}{\partial\phi^{b}}\delta\psi^{b}=-\sum_{b}\frac{\partial S}{\partial\phi^{b}}D_{A}\phi^{b}$（这里 $\delta\phi^{b}=0$ 使 $\delta(\partial S/\partial\phi^{b})=0$，且 $\partial S/\partial\phi^{b}$ 是偶的，从 $\psi$ 旁边穿过不变号）。代入 (11.12) 应得 $d\mathcal{O}_{S}^{(0)}=-\delta\big(\sum_{b}\frac{\partial S}{\partial\phi^{b}}\psi^{b}\big)$，而原文写的是 $+i\delta(\cdots)$。这个 $i$ 与 (11.8) 的系数问题同源：严格计算与原文只差一个整体因子 $i$，等价于把 $\delta$ 重新定义为 $i\delta$（注意 (11.5) 中 $\delta A=i\psi$ 已经把 $i$ 放在了 $\delta$ 的定义里）。只要在整条下降方程链中统一采用同一约定，(11.17) 的闭合性仍然成立；但若一半按 (11.5)–(11.6) 严格推导、一半照抄原文系数，符号必然打架。]

% [Tutor Note: (11.15) 第一项写作 $D_{A}\psi^{a}\wedge\psi^{b}$，而 (11.16) 又用 $\psi^{a}\wedge\psi^{b}$，两处并不同调：直接对 $\mathcal{O}_{S}^{(1)}=-\sum_{b}\frac{\partial S}{\partial\phi^{b}}\psi^{b}$ 取外微分，应为 $d\mathcal{O}_{S}^{(1)}=-\sum_{a,b}\frac{\partial^{2}S}{\partial\phi^{a}\partial\phi^{b}}D_{A}\phi^{a}\wedge\psi^{b}-\sum_{b}\frac{\partial S}{\partial\phi^{b}}D_{A}\psi^{b}$，即第一项中出现的应是 $D_{A}\phi^{a}$（再由 $\delta\psi^{a}=-D_{A}\phi^{a}$ 把它换成 $\delta\psi^{a}$），而非 $D_{A}\psi^{a}$。此外 (11.16) 中 $\frac12$ 的来历与 (11.9) 相同：$S$ 的二阶导 $\partial^{2}S/\partial\phi^{a}\partial\phi^{b}$ 关于 $a,b$ 对称，故 $\sum_{a,b}f_{ab}\delta\psi^{a}\wedge\psi^{b}=\frac12\sum_{a,b}f_{ab}\delta(\psi^{a}\wedge\psi^{b})$，因子 2 被 $\frac12$ 吸收。这两点合起来看，原文的 (11.15) 极可能是排印时把 $\phi$ 误作 $\psi$（同一页 (11.12) 的第二项也用了 $d_{A}\phi^{b}$，可资对照）；建议读者以 (11.12)、(11.16) 和 (11.17) 的闭合性为准重算一遍，并留意上一条批注所说的整体 $i$ 约定。]

这些方程可以如下概括：
\begin{equation}
  \tag{11.17}
  \label{eq:11-17}
  (d+i\delta)\left(\mathcal{O}_{S}^{(0)}+\mathcal{O}_{S}^{(1)}+\mathcal{O}_{S}^{(2)}\right)=0.
\end{equation}
（要点在于：在某个适当的双复形（double complex）中，总微分是 $d+i\delta$，而 $\mathcal{O}_{S}^{(0)}+\mathcal{O}_{S}^{(1)}+\mathcal{O}_{S}^{(2)}$ 是闭的。）

现在我们要用这一点来构造 $\mathcal{A}/\widetilde G$ 上的上同调类。
它们可以限制到平坦联络模去规范变换所得的模空间
\begin{equation}
  \tag{11.18}
  \label{eq:11-18}
  \mathcal{M}=\mathcal{A}^{\mathrm{flat}}/\widetilde G
\end{equation}
上。

若我们在 $\Sigma$ 中选取一个 $q$-圈 $C$，则由 Stokes 定理可得
\[
Q_{S}^{(q)}(C)\overset{\mathrm{def}}{=}\int_{C}\mathcal{O}_{S}^{(q)}
\]
满足
\begin{equation}
  \tag{11.19}
  \label{eq:11-19}
  \delta Q_{S}^{(q)}(C)=0,
\end{equation}
并且若 $C=\partial B$ 是一个边界，则
\begin{equation}
  \tag{11.20}
  \label{eq:11-20}
  Q_{S}^{(q)}(C)=\delta T\quad\text{对某个 }T.
\end{equation}
只要规范变换群 $\widetilde G$ 在 $\mathcal{A}$ 上的作用自由，$Q_{S}^{(q)}(C)$ 就对应于 $\mathcal{M}$ 的上同调的生成元：它们是 $\mathcal{A}/\widetilde G$ 上的上同调类，我们将把它们限制到模空间 $\mathcal{A}^{\mathrm{flat}}/\widetilde G$ 上。
这些生成元是在 Atiyah 与 Bott 的一篇奠基性论文中给出的；它们的构造方式是：取 $\mathcal{M}\times\Sigma$ 上\emph{万有丛}（universal bundle）的示性类与

% ===== 原文 p.149 =====
$\Sigma$ 的同调类作斜积（slant product）。
量子场论将计算一个生成泛函
\begin{equation}
  \tag{11.21}
  \label{eq:11-21}
  \int_{\mathcal{M}}\exp\Big\{\alpha Q_{S}^{(0)}(p)+\sum_{j=1}^{2g}\beta_{j}Q_{S'_{j}}^{(1)}(C_{j})
  +\gamma Q_{S''}^{(2)}(\Sigma)\Big\}
\end{equation}
它将编码 $\mathcal{M}$ 的上同调中所有的相交数。
这里 $p$ 是 $\Sigma$ 中的一点，$C_{j}$ 是 $\Sigma$ 中对应于同调 $H_{1}(\Sigma)$ 的 $2g$ 个圈，而 $S$、$S'_{j}$ 与 $S''$ 是 $\mathbf{g}$ 上任意的\emph{不变多项式}（它们一般是彼此不同的）。
为简单起见，我们主要处理 $G=SU(2)$ 的情形，此时 $\mathbf{g}$ 上不变多项式之环是单个生成元 $S$ 上的多项式环，$S$ 由
\[
S(\phi)=\Tr(\phi^{2})
\]
给出。

% [Tutor Note: (11.21) 把三个不同次数的类放进同一个指数里，读者容易误以为它们是同一多项式的不同分量。实际上 $Q_{S}^{(0)}(p)=S(\phi(p))$ 是零维圈（一个点）上的"积分"（即在该点取值），$Q_{S'_{j}}^{(1)}(C_{j})=\int_{C_{j}}\Tr(\phi\psi)$ 是沿 $H_{1}(\Sigma)$ 的基圈的一次积分，$Q_{S''}^{(2)}(\Sigma)$ 是 $\Sigma$ 本身（二维基本类）上的二次积分；三者分属 $H^{0}$、$H^{1}$、$H^{2}$ 的生成元。因为 $\Sigma$ 是二维的，下降序列到 $q=2$ 就截断，(11.17) 中的 $d\mathcal{O}_{S}^{(2)}=0$ 是自动的（$3$-形式在二维流形上为零），这正是"二维"这一设定的关键简化。参数 $\alpha,\beta_{j},\gamma$ 是相应的耦合常数；在 (11.27) 中作者取 $\alpha=-e^{2}/2$、$\gamma=1$，正是为了把 (11.21) 还原成 Yang–Mills 配分函数。]

\subsection{等变积分与局部化}
% 原文 §11.4：Equivariant integration and localization

现在我们回到 §11.2 的有限维情形。
我们想通过在 $M\times\mathbf{h}$ 上积分来定义一个映射 $\mathcal{I}:Z^{*}_{H}(M)\to\mathbf{C}$（其中 $Z^{*}_{H}(M)$ 是 $\Omega^{*}_{H}(M)$ 中 $D$-闭的元素）。
为了定义一个收敛的积分，我们引入收敛因子 $e^{-\varepsilon\Tr(\phi^{2})}$。
（事实上，在关于这个收敛因子的数学处理中，可以把 Gauss 型因子 $e^{-\varepsilon\Tr(\phi^{2})}$ 换成 $\mathbf{h}$ 上任一族速降函数 $\{f_{\varepsilon}\}$，只要它们在 $\varepsilon\to0$ 时表示 Dirac delta 分布。）
对 $\alpha\in Z^{*}_{H}(M)$ 我们定义
\begin{equation}
  \tag{11.22}
  \label{eq:11-22}
  \mathcal{I}(\alpha)=\int_{\phi\in\mathbf{h}}\int_{M}d\phi_{1}\dots d\phi_{s}\,
  e^{-\varepsilon\Tr(\phi^{2})}\alpha(\phi).
\end{equation}

事实上，为了保证积分 (11.22) 收敛，我们必须对 $\alpha$ 加一些假设。
我们将允许 $\alpha$ 在 $\phi$ 很大时关于 $\phi$ 至多指数增长。

% [Tutor Note: $\mathbf{h}$ 的维数记号在此处前后不一：(11.22) 写作 $d\phi_{1}\dots d\phi_{s}$，而同一测度在 (11.25)、(11.26) 中写作 $d\phi_{1}\dots d\phi_{n}$；按上下文 $s$ 与 $n$ 都应理解为 $\dim\mathbf{h}$，转写时保留原样，读者不要误以为出现了两种不同的积分变量。另需注意，这个"积分"只是无穷维路径积分的有限维模型：真正的 $\mathbf{h}$ 是 $\mathrm{Lie}(\widetilde G)$ 的类比，其维数在 $\dim\Sigma\to\infty$ 的意义下发散，故 $\mathcal{I}(\alpha)$ 只能作为形式的（由 (11.25) 的 $\delta(\mu)$ 修正的）定义；作者随后正是用 $e^{-\varepsilon\Tr(\phi^{2})}$ 这一个"Gauss 截断"来替代对 $\dim\mathbf{h}$ 的正规化。]

一类有用的等变闭形式来自下述特殊情形：$M$ 是带辛形式 $\omega$ 的辛流形（symplectic manifold），并且 $H$ 在 $M$ 上的作用是 Hamilton 作用，矩映射（moment map）为 $\mu:M\to\mathbf{h}^{*}$。我们容易看出
\begin{equation}
  \tag{11.23}
  \label{eq:11-23}
  \widetilde\omega=-i\Tr(\phi\mu)+\omega
\end{equation}
满足
\[
D\widetilde\omega=0
\]
（其中 $D$ 是由 (11.10) 定义的 Cartan 模型微分）。
在辛流形上，我们有时取 $\alpha$ 为 $\widetilde\omega$ 的指数乘以一个关于 $\phi$ 为多项式的形式。

积分 (11.22) 的计算将由一个局部化原理支配，它类似于导致辛流形上带 Hamilton 环面作用的振荡积分的 Duistermaat–Heckman 公式的那个原理。
我们选取一个 $H$-不变的 1-形式 $\lambda$。对这种形式，
\begin{equation}
  \tag{11.24}
  \label{eq:11-24}
  D\lambda=d\lambda-i(V(\phi),\lambda);
\end{equation}
其中 $V(\phi)$ 是由 $\phi$ 的作用给出的向量场（它当然线性依赖于 $\phi$）。
把 $\exp(-\varepsilon\Tr(\phi^{2}))$ 换成 $\exp(-\varepsilon\Tr(\phi^{2})+itD\lambda)$ 并不改变积分 (11.22)，因为 $\exp(itD\lambda)$ 展开为 $1+D(it\lambda+\dots)$。
由于 $-tD\lambda=it(V(\phi),\lambda)-t(d\lambda)$，当 $t\to\infty$ 时积分具有剧烈振荡的相位，于是它局部化到满足 $(V(\phi),\lambda)=0$ 的点附近。
可以直接由驻相（stationary phase）论证看出这种局部化。
或者，也可以先做 $\phi$ 积分

% ===== 原文 p.150 =====
——利用指数关于 $\phi$ 是二次的这一事实。
配方并完成对 $\phi$ 的积分之后，我们得到一个因子
\[
\exp\left(-\frac{t^{2}}{4\varepsilon}\sum_{a}(V_{a},\lambda)^{2}\right)
\]
（以及其它对 $t$ 的依赖弱得多的因子），于是考察 $t\to\infty$ 的行为，我们看到积分局部化到使所有 $a$ 都满足 $(V_{a},\lambda)=0$ 的点附近。

% [Tutor Note: 原文此处连写"and and"（p.150 第 2 行 "completing the square and and performing the integral over $\phi$"），系排印重复，译文未保留。更值得注意的是配方结果中 $t^{2}$ 项的符号：$\phi$ 的二次部分由 $-\varepsilon\Tr(\phi^{2})$ 与 $t(V(\phi),\lambda)=t\sum_{a}\phi^{a}(V_{a},\lambda)$ 构成，配方给出 $-\varepsilon\big(\phi^{a}-\frac{t}{2\varepsilon}(V_{a},\lambda)\big)^{2}$ 再加上 $\pm\frac{t^{2}}{4\varepsilon}\sum_{a}(V_{a},\lambda)^{2}$，符号取决于交叉项是实是虚，也就是取决于 (11.10) 中的 $i$ 是否保留（见 p.147 脚注 1）以及 $\phi$ 的围道取实轴还是虚轴。原文写出的带负号的 $\exp\big(-\frac{t^{2}}{4\varepsilon}\sum_{a}(V_{a},\lambda)^{2}\big)$ 要求交叉项为纯虚，即采用数学家略去 $i$ 的约定 $D=d+\iota_{V(\phi)}$；若严格按 (11.10)、(11.24) 的 $D=d-i\iota$ 直接计算，得到的是带正号的指数，此时必须把 $\phi$ 的积分围道旋到虚轴（等价于对 $-\varepsilon\Tr(\phi^{2})$ 作解析延拓）才能恢复衰减。做局部化计算前务必先固定这一约定，否则会得到"指数增长"而非"局部化"的荒谬结论。]

现在我们专门化到下述情形：$M$ 是辛流形，配有紧群 $H$ 的 Hamilton 作用，矩映射记作 $\mu$。
我们定义 $I=|\mu|^{2}:M\to\mathbf{R}$，即矩映射的模平方。
我们可以选取一个与辛结构相容的殆复结构（almost complex structure）$J\in\mathrm{End}(TM)$（所谓相容，指辛形式关于 $J$ 是 $(1,1)$ 型的并且是正的），并令 $\lambda=J(dI)$ 为上面定义的不变 1-形式。
使 $(V_{a},\lambda)=0$ 的点恰好可以等同于 $I$ 的临界点。
我们看出这一点，是因为
\[
dI=2\sum_{a}\mu_{a}\,d\mu_{a},
\]
于是由 $V_{a}=\omega^{-1}d\mu_{a}$（这里把辛形式 $\omega$ 看作对每个 $x\in M$ 的同构 $T_{x}M\to T^{*}_{x}M$，其逆为 $\omega^{-1}:T^{*}_{x}M\to T_{x}M$）可得
\[
\omega^{-1}dI=2\sum_{a}\mu_{a}V_{a}.
\]
这里 $dI\in\Omega^{1}(M)$，故应把 $\omega^{-1}dI$ 看作 $\Gamma(TM)$ 的元素。
于是 $(V_{a},\lambda)=0$ 蕴涵 $(\omega^{-1}dI,\lambda)=0$，换言之
\[
(\omega^{-1}dI,J(dI))=0
\]
（此式两边都在 $\Gamma(TM)$ 中）。
由于 $\omega$ 与 $J$ 合起来在余切空间上给出一个非退化对称双线性形式 $g$，上式蕴涵 $dI=0$；$g$ 由
\[
g(X,Y)=\omega^{-1}(X,JY)
\]
定义（对 $X,Y\in\Gamma(T^{*}M)$）。
因此我们的积分在形式上局部化到 $I$ 的临界点上。
这类临界点的一个明显来源是 $\mu^{-1}(0)$，但也存在 $I\neq0$ 的其它临界点。

现在我们来给出 Duistermaat–Heckman 公式（的一个版本）的推导。
如上所述，形式 $\widetilde\omega=-i\Tr(\phi\mu)+\omega$（已在 (11.23) 中引入）是等变闭的。
因此我们可以计算
\begin{equation}
  \tag{11.25}
  \label{eq:11-25}
  \frac{1}{\mathrm{Vol}(H)}\int_{\phi\in\mathbf{h}}d\phi_{1}\dots d\phi_{n}\int_{M}
  \exp\left(-i\Tr(\phi\mu)+\omega\right).
\end{equation}
（注意：若在公式 (11.25) 中不包含对 $\phi$ 的积分，则 $M$ 上的振荡积分将对应于对 $T$ 的不动点集的各连通分支求和：可参见参考文献中所列 Berline、Getzler 与 Vergne 的书第 7 章。）
然而，如果我们先对 $\phi$ 积分，就得到一个 delta 函数 $\delta(\mu)$，于是公式成为
\[
\frac{1}{\mathrm{Vol}(H)}\int_{M}(\exp\omega)\,\delta(\mu)=\int_{\mu^{-1}(0)/H}e^{\omega}
\]

% ===== 原文 p.151 =====
换言之，即辛商（symplectic quotient）$\mu^{-1}(0)/H$ 的辛（Liouville）体积。

然而我们还希望在积分核中纳入一个因子 $\exp(-\varepsilon\Tr(\phi^{2}))$。
因此我们希望计算
\begin{equation}
  \tag{11.26}
  \label{eq:11-26}
  Z(\varepsilon)\overset{\mathrm{def}}{=}\frac{1}{\mathrm{Vol}(H)}
  \int_{\phi\in\mathbf{h}}d\phi_{1}\dots d\phi_{n}\int_{M}\alpha\,
  \exp\left(-i\Tr(\phi\mu)+\omega-\varepsilon\Tr(\phi^{2})\right),
\end{equation}
其中 $\alpha=\alpha(\phi)$ 是 $M$ 上的等变闭形式，它对 $\phi$ 有\emph{多项式}依赖。
我们假定 $H$ 在 $\mu^{-1}(0)$ 上自由作用；在 $\mu^{-1}(0)$ 附近，等变闭类 $\Tr(\phi^{2})$（它来自 Cartan 模型中一个点的 $H$-等变上同调）是辛商 $\mu^{-1}(0)/H$ 上某个类 $\Theta$ 的拉回，而这个类在 Chern–Weil 理论中会被计算为 $\Tr(F_{B}^{2})$，其中 $B$ 是主 $H$-丛 $\mu^{-1}(0)\to\mu^{-1}(0)/H$ 上的联络，$F_{B}$ 是它的曲率。
类似地，在 $\mu^{-1}(0)$ 附近，$\alpha$ 限制为 $\mu^{-1}(0)/H$ 上某个类 $\widehat\alpha$ 的拉回。

我们发现
\[
Z(\varepsilon)=\int_{\mu^{-1}(0)/H}\widehat\alpha\,\exp(\omega-\varepsilon\Theta)
+O(e^{-b/\varepsilon}),
\]
其中 $b$ 是 $I$ 的最小非零临界值。
为了启发式地看出为什么 $I$ 的非零临界值 $b$ 应当对积分 (11.26) 给出量级为 $e^{-b/\varepsilon}$ 的贡献，我们通过先积掉 $\phi$ 来计算 (11.26)（为简单起见设 $\alpha=1$）。
我们用在 $\phi$ 中配方的方法来计算关于 $\phi$ 的积分，得到
\[
Z(\varepsilon)=\int_{M}\exp\left(\omega-\frac{1}{\varepsilon}\Tr(\mu^{2})\right).
\]
于是，粗略地说，$\mu$ 非零的临界点所贡献的量在 $\varepsilon$ 很小时应为 $\exp\left(-\Tr(\mu^{2})/\varepsilon\right)$ 的量级。

\subsection{等变积分：无穷维情形}
% 原文 §11.5：Equivariant integration: the infinite dimensional case

现在我们过渡到无穷维情形，研究对应于等变积分的路径积分。
回想我们已定义生成泛函 (11.21)。
我们限制到 $G=SU(2)$ 或 $SO(3)$；我们将取不变多项式 $S$、$S'_{j}$ 与 $S''$ 全都等于生成多项式（记作 $S$），即 $S(\phi)=\Tr(\phi^{2})$。
（我们已在 $\mathbf{g}$ 上取定一个 $\mathrm{Ad}$-不变内积，记作 $\Tr$；$G$ 上相应地定义了测度。）
我们假定 (11.21) 中 $\gamma=1$、$\alpha=-e^{2}/2$，并注意到对 $S=\Tr(\phi^{2})$ 有
\[
Q_{S}^{(1)}(C_{j})=\int_{C_{j}}\Tr(\phi\psi).
\]
于是给出 (11.21) 的路径积分成为（这里 $p$ 是 $\Sigma$ 中的一点）
\begin{equation}
  \tag{11.27}
  \label{eq:11-27}
  \begin{aligned}
  Z(e,\{\beta_{j}\})={}&\frac{1}{\mathrm{Vol}(\widetilde G)}
  \int\mathcal{D}A\,\mathcal{D}\psi\,\mathcal{D}\phi\,
  \exp\bigg(-\frac{e^{2}}{2}\Tr(\phi^{2})(p)\\
  &\qquad+\sum_{j=1}^{2g}\beta_{j}\int_{C_{j}}\Tr(\phi\psi)+\int_{\Sigma}Q_{S}^{(2)}\bigg).
  \end{aligned}
\end{equation}
（注意：当 $S(\phi)=\Tr(\phi^{2})$ 时，对象 $\int_{\Sigma}Q_{S}^{(2)}$ 对应于 $\mathcal{A}$ 上的辛形式，并将等同于模空间 $\mathcal{M}$ 上的标准辛形式。）

% [Tutor Note: (11.27) 的第一项是 $-\frac{e^{2}}{2}\Tr(\phi^{2})(p)$，看似缺少积分号，实则不然：$Q_{S}^{(0)}(p)=S(\phi(p))$ 是定义在零维圈（一个点 $p$）上的 $0$-形式，故"积分"就是取值。把它与 (11.3) 对照可见二者的差别：在 (11.3) 中该项是 $-\frac{e^{2}}{2}\int d\mu\Tr(\phi^{2})$（对整个曲面积分），而这里只在一个点取值。这一替换之所以正当，是因为在 (11.21) 的生成泛函中 $Q_{S}^{(0)}$ 只以 $\alpha Q_{S}^{(0)}(p)$ 的形式出现，而由下降方程 $Q_{S}^{(0)}$ 与点的选取只差一个 $\delta$-恰当项。但要注意：正因为只取一点的值，$\phi^{2}$ 的"收敛因子"在二维场论中不是良定义的（$\phi$ 的快度发散），所以严格说来 (11.27) 必须理解为后面 (11.28) 中 $\varepsilon$ 位移的简写，而不能逐点理解。]

% ===== 原文 p.152 =====
我们现在可以看出，积分 (11.27) 如何能被改写成这样的形式：使涉及 $q=0,1,2$ 的 $Q_{S}^{(q)}$ 的生成泛函 (11.21) 等于一个只涉及 $q=0,2$ 的 $Q_{S}^{(q)}$ 的相应生成泛函。
这是通过对 $\psi$ 作积分来实现的，而该积分是 Gauss 型的：通过配方完成这个积分之后我们看到，我们得到
\begin{equation}
  \tag{11.28}
  \label{eq:11-28}
  \frac{1}{\mathrm{Vol}(\widetilde G)}\int\mathcal{D}A\,\mathcal{D}\phi\,
  \exp\left\{i\int_{\Sigma}\Tr(\phi F_{A})-\varepsilon\Tr(\phi^{2})(p)\right\},
\end{equation}
其中 $\varepsilon=e^{2}/2-\sum_{i<j}\beta_{i}\beta_{j}\#(C_{i}\cap C_{j})$。
换言之，我们可以消去模空间上同调环中对应于 $Q_{S}^{(1)}$ 的奇维生成元（即模空间中的非代数圈）。
这一结果可以用一个简单的拓扑论证来解释。

若推广到秩大于 1 的 Lie 群的情形，人们仍然可以通过对 $\psi$ 作 Gauss 积分，把奇维圈消去而代之以偶维圈；但在这种情形下，对这一操作的拓扑解释尚不清楚。

% [Tutor Note: $\varepsilon$ 的位移公式 $\varepsilon=\frac{e^{2}}{2}-\sum_{i<j}\beta_{i}\beta_{j}\#(C_{i}\cap C_{j})$ 是 p.152 最需要留意的地方，其来源如下：把 (11.27) 中 $\exp\big(\sum_{j}\beta_{j}\int_{C_{j}}\Tr(\phi\psi)\big)$ 展开，二次项 $\frac12\sum_{i,j}\beta_{i}\beta_{j}\int_{C_{i}}\Tr(\phi\psi)\int_{C_{j}}\Tr(\phi\psi)$ 与 $\psi$ 的 Gauss 积分（(11.4) 中的 $\frac12\int\Tr(\psi\wedge\psi)$ 提供传播子 $\langle\psi\psi\rangle\sim\delta^{(2)}$）收缩，得到 $\sum_{i,j}\beta_{i}\beta_{j}\int_{C_{i}}\!\!\int_{C_{j}}\delta^{(2)}$。当 $i\neq j$ 且两圈横截相交时，这个双重积分恰为相交数 $\#(C_{i}\cap C_{j})$；对 $i<j$ 成对求和即得上式。换言之，正是"$\psi$ 的传播子是 $\Sigma$ 上的 delta 函数"这一事实，使得两个一维圈上的算子收缩出零维的交点，从而把 $Q^{(1)}$ 的偶次幂并入 $Q^{(0)}$ 的系数 $\varepsilon$ 中。这也解释了作者随后所说的"非代数圈被消去"：在 $G=SU(2)$ 时 $H^{1}$ 类由 $2g$ 个基本圈生成，而它们的乘积恰由交形式（intersection form）$\#(C_{i}\cap C_{j})$ 完全刻画，故只留下偶数维的代数圈。]

\subsection{Yang–Mills 理论的配分函数}
% 原文 §11.6：The partition function of Yang-Mills theory

现在我们已经把需要计算的公式化归为 Yang–Mills 理论（不含费米型变量 $\psi$）的配分函数 (11.2)。
为了计算配分函数，我们在 Hamilton 框架中量子化这个理论，然后把各部分拼起来以得到整体结果。

为了引入 Hamilton 框架，我们把 Riemann 曲面换成一个柱面 $C\times[0,T]$（其中 $C\cong S^{1}$ 应被看作空间，并假定其周长为 $L$，而 $[0,T]$ 是长度为 $T$ 的时间区间）。
固定时刻的场空间等于圆 $C$ 上联络的空间 $\mathcal{A}_{C}$，故 Hilbert 空间是
\[
\mathcal{H}=C^{\infty}(\mathcal{A}_{C})^{\widetilde G},
\]
即 $\mathcal{A}_{C}$ 上的规范不变函数。
由于在规范等价意义下一切联络都由其绕曲线 $C$ 的和乐（holonomy）分类，我们发现 Hilbert 空间等同于
\[
\mathcal{H}=C^{\infty}(G)^{G},
\]
即 $G$ 上在伴随作用下不变的那些函数构成的环（等价地，即极大环面 $T$ 上在 Weyl 群作用下不变的那些函数）。
于是 Hilbert 空间等同于 $G$ 的不可约（复）表示 $R$（或等价地等同于它们的特征标，记作 $\chi_{R}$ 或 $\Tr_{R}$）：我们取特征标 $\chi_{R}$ 构成 $\mathcal{H}$ 的一组基，即
\[
\mathcal{H}=\oplus_{R}\mathbf{C}_{(R)}.
\]
对每个不可约表示 $R$，基元 $\chi_{R}$ 等同于 $\mathcal{A}_{C}$ 上的一个函数 $\Psi_{R}$，即
\[
\Psi_{R}(A)=\chi_{R}(\mathrm{Hol}_{C}A).
\]
理论的 Hamilton 量为
\[
H=-\frac{e^{2}}{2}\int_{C}\Tr|F_{01}|^{2}
\]

% ===== 原文 p.153 =====
（其中 $F_{01}$ 是联络的曲率），它化为（把 Hamilton 量看作 $\mathcal{H}$ 上的算子，用 $\mathcal{A}_{C}$ 上的位置变量 $A^{b}(\theta)d\theta$ 以及相应的动量变量 $(\partial A^{b}(\theta)/\partial t)d\theta$ 表示，而后者在正则量子化下被等同于 $-i\partial/\partial A^{b}(\theta)$）
\[
H=\frac{e^{2}}{2}\int_{C}\Tr\left(\frac{\delta}{\delta A}\right)^{2}.
\]
算子 $\delta/\delta A^{a}(\theta)$ 作用在 $\Psi_{R}$ 上给出
\[
\frac{\delta}{\delta A^{a}(\theta)}\Tr_{R}\mathrm{Hol}_{C}(A)
=\Tr_{R}T^{a}\mathrm{Hol}_{C}(A),
\]
这里用到了 $G$ 的李代数 $\mathbf{g}$ 的一组基 $T^{a}$。我们因此看到
\begin{equation}
  \tag{11.29}
  \label{eq:11-29}
  \sum_{a}\frac{\delta}{\delta A^{a}(\theta)}\frac{\delta}{\delta A^{a}(\theta)}
  \Tr_{R}\mathrm{Hol}_{C}(A)=\Tr_{R}\Big(\sum_{a}T^{a}T^{a}\mathrm{Hol}_{C}(A)\Big).
\end{equation}
（这实际上是第 5 讲的一个结果，那里考虑过同一个问题。）
元素 $\sum_{a}T^{a}T^{a}$ 是 $G$ 的\emph{二次 Casimir 算符}（quadratic Casimir），它在表示 $R$ 中的取值记作
\[
\sum_{a}T^{a}T^{a}\overset{\mathrm{def}}{=}c_{2}(R).
\]
于是我们看到
\begin{equation}
  \tag{11.30}
  \label{eq:11-30}
  H\Psi_{R}=\frac{e^{2}c_{2}(R)L}{2}\Psi_{R}.
\end{equation}
我们因此看到
\begin{equation}
  \tag{11.31}
  \label{eq:11-31}
  \langle\Psi_{R}|e^{-HT}|\Psi_{R}\rangle
  =\exp\left(-e^{2}c_{2}(R)LT/2\right)=\exp\left(-e^{2}\mathbf{a}c_{2}(R)/2\right),
\end{equation}
其中 $\mathbf{a}=LT$ 是柱面的面积。答案只依赖于组合 $e^{2}\mathbf{a}$。

% [Tutor Note: 从 $H=-\frac{e^{2}}{2}\int_{C}\Tr|F_{01}|^{2}$ 到 $H=+\frac{e^{2}}{2}\int_{C}\Tr(\delta/\delta A)^{2}$，符号翻转是正则量子化的直接结果，也是本页最容易出错的地方：$F_{01}=\partial_{0}A_{1}-\partial_{1}A_{0}+[A_{0},A_{1}]$ 在 $A_{0}=0$ 规范下等于 $\partial_{t}A$，而动量 $\partial_{t}A^{b}$ 被等同于 $-i\delta/\delta A^{b}$，故 $(F_{01})^{2}\mapsto(-i)^{2}(\delta/\delta A)^{2}=-(\delta/\delta A)^{2}$，与前面已有的负号相消。若忘记 $(-i)^{2}=-1$，就会得到 $H\Psi_{R}=-\frac{e^{2}c_{2}(R)L}{2}\Psi_{R}$，从而 (11.31) 变成指数增长——这与"二维 Yang–Mills 禁闭、Wilson 圈呈面积定律衰减"的物理图像直接矛盾，可作为自检的判据。另外 (11.31) 中的面积用了粗体 $\mathbf{a}$，勿与表示指标 $a$、圈指标 $i,j$ 混淆；原文在 p.145 与 p.153 都用粗体 $\mathbf{a}$ 表示面积。]

尽管我们今天的兴趣主要在拓扑方面，我们仍要停下来指出方程 (11.30) 的一个物理解释。
它断言：若 $R$ 是 $G$ 的非平凡表示，则态 $\Psi_{R}$ 的能量正比于 $L$（当 $L\to\infty$ 时，事实上精确正比于 $L$）。
这等价于说二维的纯 Yang–Mills 理论是禁闭的（confining）；它对所有表示 $R$ 都表现出面积定律（area law）。
（二维 Yang–Mills 理论不描述任何传播的自由度，因此不存在由理论中的电荷对外部电荷进行屏蔽的可能；这就是为什么要得到这个结果，我们不必限制到"$R$ 是 $G$ 的某个覆盖群的表示、而不是 $G$ 本身的表示"这一情形。当存在物质场时为何需要这一限制，已在第 7 讲末尾解释过。）

事实上，这个量子场论中重整化的效应是在积分核中引入形如
\begin{equation}
  \tag{11.32}
  \label{eq:11-32}
  \exp\left(\alpha\int_{\Sigma}\mathcal{R}\,d\mu+(e^{2}/2)\beta\int_{\Sigma}d\mu\right)
\end{equation}
的因子（其中 $\mathcal{R}$ 是 Riemann 曲率，$\alpha$ 与 $\beta$ 是适当选取的实值常数）。
这个理论是超可重整的（superrenormalizable），故重整化只会引入比经典 Lagrange 量中已有项维数更低的项；我们刚刚写下的这些项是唯一可能的项。
含 $\beta$ 的项的作用是把二次 Casimir $c_{R}$ 改变为一个重整化值 $\widetilde c_{R}$，它与 $c_{R}$ 相差一个常数。
因子 $\alpha$ 乘的是 $\Sigma$ 的 Euler 示性数：它最终将被调整，以确保整体的

% [原文此句在 p.153 末尾中断，续于 p.154（本任务范围之外）。]
